\chapter{Sources, préparation et traçabilité}
\section{Séries mensuelles d'IPC}
Les trois classeurs mensuels du HCP couvrent janvier 2017 à octobre 2026 \cite{hcpIndicateurs}. Les grandes divisions fournissent l'indice général, l'alimentation et les produits non alimentaires. Les deux groupes de divisions fournissent six séries chacun. Le fichier annuel de détail « Biens et services divers » éclaire la nomenclature. Il ne complète pas la série mensuelle.

Le script \texttt{clean\_ipc.py} lit les en-têtes à la ligne 24 et les observations à partir de la ligne 25. Il élimine les cellules de présentation et la note de source, transforme le mois \texttt{AAAA/MM} en date au premier du mois et trie la série par ordre croissant. Il contrôle les intitulés, les doublons de dates, la présence des 118 mois attendus et les valeurs numériques positives et finies. Une fusion sur la date associe les 118 lignes de chaque tableau. Les libellés HCP sont conservés, y compris la distinction entre l'agrégat « Alimentation » et sa division « Produits alimentaires et boissons non alcoolisées ». Le CSV final compte 118 lignes, 15 séries et une colonne date. Il ne comporte ni mois dupliqué, ni valeur manquante. Les 1\,770 indices ont été rapprochés des cellules Excel lors de la préparation.

\section{Enquête sur le niveau de vie des ménages}
L'ENNVM a été collectée du 15 mars 2022 au 14 mars 2023 \cite{hcpENNVM}. Le classeur local contient les tabulations agrégées. La nomenclature COICOP Maroc 2022 organise les postes par fonction de consommation. La dépense annuelle moyenne par personne (DAMP) est exprimée en dirhams courants par personne et par an. Le coefficient budgétaire (CB) représente la part d'un poste dans le budget publié. Les valeurs constituent une coupe agrégée sur la période de l'enquête.

Sur les 24 feuilles inspectées, le tableau SD~12 de la feuille « Partie3-Structure de la DAMP » fournit les groupes de dépenses par quintile et milieu. Il commence à la ligne 288. Les blocs de données occupent les lignes 294--308, 315--329 et 336--350. Le script \texttt{clean\_households.py} vérifie les trois agrégats avec SD~01, qui publie les mêmes coefficients et arrondit davantage la DAMP. La feuille « Considérations importantes » donne les seuils de quintile pour chaque périmètre. Les feuilles DAMPG01 à DAMPG14 détaillent les catégories utilisées pour examiner la correspondance.

Les 18 profils comprennent un ensemble et cinq quintiles dans chacun des périmètres national, urbain et rural. Les quintiles sont définis selon la \emph{dépense annuelle par personne}. Chaque classement utilise des seuils propres à son périmètre. La comparaison de C1 urbain à C1 rural porte ainsi sur une position relative au sein de deux distributions. Les 18 profils ne sont pas 18 observations individuelles, et leurs ensembles et quintiles se recouvrent en population. Les 14 coefficients de chaque profil totalisent de 99,8 à 100,1\,\%. Sous l'hypothèse d'un arrondi au dixième, une borne maximale de $14\times0,05=0,70$ point s'applique à l'écart de somme. Les coefficients publiés ont été conservés.

\section{Chaîne d'analyse}
Le notebook 01 sert de trame exploratoire. Le notebook 02 calcule les statistiques à partir du CSV IPC. Le notebook 03 documente la couverture des profils et l'audit des correspondances. Le notebook 04 calcule les trajectoires historiques et les chocs hypothétiques. L'annexe \ref{ann:trace} relie les tableaux et figures aux fichiers du projet.
